\section{Term calculus (CoC + inductives)}
\label{sec:calculus}

This section presents the term calculus used throughout the paper: a Calculus of Constructions (CoC) extended with inductive types and fixpoints. This is the source language for both the natural deduction typing system and the sequent calculus developed in Section~\ref{sec:focused}.

\subsection{Syntax}

Terms are defined inductively as follows:
\[
\begin{array}{rcl}
t, A, B &::=& x \mid \Sort{i} \mid \PiTy{A}{B} \mid \Lam{A}{t} \mid \App{t}{u} \\
        &\mid& \Fix{A}{t} \mid \Ind{I}{\vec{a}} \mid \Roll{I}{c}{\vec{a}} \mid \Case{I}{t}{C}{\vec{b}}
\end{array}
\]

\noindent where:
\begin{itemize}[leftmargin=2em]
\item $x$ ranges over de Bruijn indices (variables)
\item $\Sort{i}$ is the universe at level $i$ (we use a cumulative hierarchy)
\item $\PiTy{A}{B}$ is the dependent function type with domain $A$ and codomain $B$ (binding one variable)
\item $\Lam{A}{t}$ is a lambda abstraction with explicit type annotation $A$
\item $\App{t}{u}$ is function application
\item $\Fix{A}{t}$ is a fixpoint with type annotation $A$ (binding one variable for recursion)
\item $\Ind{I}{\vec{a}}$ is an applied inductive type $I$ with parameters/indices $\vec{a}$
\item $\Roll{I}{c}{\vec{a}}$ is a constructor application (constructor $c$ of inductive $I$)
\item $\Case{I}{t}{C}{\vec{b}}$ is dependent pattern matching on scrutinee $t$ with motive $C$ and branches $\vec{b}$
\end{itemize}

\noindent\textbf{Conventions:}
We write $t[u/x]$ for capture-avoiding substitution and $t[u/]$ for substitution under one binder (replacing de Bruijn index 0).
Contexts $\Gamma$ are lists of types. The global environment $\Sigma$ records inductive definitions (constructors, elimination rules, strict positivity).

\subsection{Natural deduction typing rules}

The typing judgement $\hastype{\Sigma}{\Gamma}{t}{A}$ states that term $t$ has type $A$ in context $\Gamma$ under global environment $\Sigma$. We abbreviate this as $\Gamma \vdash t : A$ when $\Sigma$ is clear from context.

\subsubsection{Structural rules}

\begin{Rules}
\Infer{Var}{
  x : A \in \Gamma
}{
  \Gamma \vdash x : A
}
\\[1.5em]
\Infer{Weak}{
  \Gamma \vdash t : A \quad
  \Gamma \vdash B : \Sort{i}
}{
  \Gamma, B \vdash t : A
}
\end{Rules}

\subsubsection{Universe and conversion}

\begin{Rules}
\Infer{Sort}{
  i < j
}{
  \Gamma \vdash \Sort{i} : \Sort{j}
}
\\[1.5em]
\Infer{Conv}{
  \Gamma \vdash t : A \quad
  \Gamma \vdash B : \Sort{i} \quad
  A \equiv_\beta B
}{
  \Gamma \vdash t : B
}
\end{Rules}

\noindent Here $A \equiv_\beta B$ denotes $\beta$-convertibility (the reflexive-symmetric-transitive closure of $\beta$-reduction).

\subsubsection{Dependent functions ($\Pi$-types)}

\begin{Rules}
\Infer{Pi-Form}{
  \Gamma \vdash A : \Sort{i} \quad
  \Gamma, A \vdash B : \Sort{j}
}{
  \Gamma \vdash \PiTy{A}{B} : \Sort{\max(i,j)}
}
\\[1.5em]
\Infer{Pi-Intro}{
  \Gamma \vdash \PiTy{A}{B} : \Sort{i} \quad
  \Gamma, A \vdash t : B
}{
  \Gamma \vdash \Lam{A}{t} : \PiTy{A}{B}
}
\\[1.5em]
\Infer{Pi-Elim}{
  \Gamma \vdash t : \PiTy{A}{B} \quad
  \Gamma \vdash u : A
}{
  \Gamma \vdash \App{t}{u} : B[u/]
}
\end{Rules}

\subsubsection{Fixpoints (recursion)}

\begin{Rules}
\Infer{Fix-Form}{
  \Gamma \vdash A : \Sort{i} \quad
  \Gamma, A \vdash t : A
}{
  \Gamma \vdash \Fix{A}{t} : A
}
\end{Rules}

\noindent The fixpoint $\Fix{A}{t}$ represents a recursive definition where the recursive call is represented by de Bruijn index 0 under the binder.

\subsubsection{Inductive types}

Let $\Sigma$ define an inductive type $I$ with constructors $\{c_0, \ldots, c_{n-1}\}$. For each constructor $c$, we have a type schema $T_c$ recorded in $\Sigma$.

\begin{Rules}
\Infer{Ind-Form}{
  \Sigma(I) = (\vec{p}, \vec{idx}, \ldots) \quad
  \Gamma \vdash \vec{a} : \vec{\tau}
}{
  \Gamma \vdash \Ind{I}{\vec{a}} : \Sort{i}
}
\\[1.5em]
\Infer{Ind-Intro}{
  \Sigma(I) = (\ldots, c : T_c, \ldots) \quad
  \Gamma \vdash \vec{a} : T_c[\vec{p}/]
}{
  \Gamma \vdash \Roll{I}{c}{\vec{a}} : \Ind{I}{\vec{p}}
}
\end{Rules}

\subsubsection{Pattern matching (elimination)}

\begin{Rules}
\Infer{Ind-Elim}{
  \begin{array}{l}
  \Gamma \vdash t : \Ind{I}{\vec{a}} \\
  \Gamma, \Ind{I}{\vec{a}} \vdash C : \Sort{j} \\
  \forall c \in \text{constrs}(I),\; \Gamma \vdash b_c : T_c \to C[\Roll{I}{c}{\cdots}/]
  \end{array}
}{
  \Gamma \vdash \Case{I}{t}{C}{\vec{b}} : C[t/]
}
\end{Rules}

\noindent The motive $C$ is a type former depending on the scrutinee. Each branch $b_c$ handles constructor $c$, and the overall case expression has type $C[t/]$ (the motive instantiated at the scrutinee).

\subsection{Operational semantics}

We use call-by-name (CBN) reduction, which is essential for the supercompilation correspondence (driving = lazy unfolding).

\subsubsection{Reduction rules}

\[
\begin{array}{rcl}
\App{(\Lam{A}{t})}{u} &\to& t[u/] \qquad \text{($\beta$-reduction)} \\[0.5em]
\Fix{A}{t} &\to& t[\Fix{A}{t}/] \qquad \text{(fixpoint unfolding)} \\[0.5em]
\Case{I}{(\Roll{I}{c}{\vec{a}})}{C}{\vec{b}} &\to& b_c\,\vec{a} \qquad \text{($\iota$-reduction)}
\end{array}
\]

\noindent CBN reduces at the outermost (head) position: we contract a $\beta$-redex only when the function is a $\lambda$, we unfold a fixpoint only at the head, and we perform $\iota$-reduction only when the scrutinee is a constructor.
We do \emph{not} reduce under binders in the operational semantics.
(Independently, the typing rules use the standard conversion rule based on $\beta$-convertibility.)

\subsection{Observational equivalence}

Two terms are observationally equivalent if they produce indistinguishable results in all closing contexts. We use the CIU (Closed Instantiations of Uses) theorem from the mechanization: terms are equivalent iff they pass the same observations at the same types.

For this paper, the key observation predicate is:
\[
\mathsf{Obs}_{\Nat}(t) \iff t \steps \Roll{\Nat}{\mathsf{zero}}{} \;\vee\; t \steps \Roll{\Nat}{\mathsf{succ}}{n} \text{ for some } n
\]

\noindent Supercompilation preserves observational equivalence: the residual program produced by the supercompiler is observationally equivalent to the input.
